\begin{center}
    \textbf{Masterarbeit} \\
    \vspace{11pt}
    für Herrn Max Frankenberg (Matrikel-Nr. 2419961) \\
    \vspace{11pt}
    \textbf{Data-Driven Model Predictive Control for Seizure Suppression with Sparse Stimulation} \\
\end{center}

% Aufgabenstellung:

Fokale epileptische Anfälle entstehen in einem begrenzten Bereich des Gehirns und können sich im weiteren Verlauf auf verbundene Hirnregionen ausbreiten. 
Die gezielte elektrische Stimulation bietet einen möglichen Ansatz, um diese Ausbreitung zu unterdrücken und damit die Auswirkungen epileptischer Anfälle zu begrenzen. 
Hierfür sind geschlossene Regelungskonzepte erforderlich, die anhand gemessener neuronaler Aktivität geeignete Stimulationssignale in Echtzeit bestimmen. 
Eine besondere Herausforderung besteht darin, dass für die Regelung lediglich indirekte Messungen der Hirnaktivität, beispielsweise mittels Elektroenzephalographie (EEG), zur Verfügung stehen. 
Modellprädiktive Regler (Model Predictive Control, MPC) ermöglichen es, zukünftige Systementwicklungen vorherzusagen und die Stimulation unter Berücksichtigung von Nebenbedingungen gezielt anzupassen. 
Voraussetzung hierfür sind geeignete Vorhersagemodelle, die den Zusammenhang zwischen der gemessenen Hirnaktivität und den applizierten Stimulationssignalen beschreiben.

Im Rahmen dieser Masterarbeit soll der Einsatz modellprädiktiver Regelung zur Unterdrückung der Ausbreitung fokaler epileptischer Anfälle untersucht werden. 
Hierzu sollen zunächst datengetriebene Vorhersagemodelle für EEG-Signale und deren spektrale Eigenschaften entwickelt und miteinander verglichen werden. 
Darauf aufbauend sollen MPC-Ansätze zur Bestimmung geeigneter Stimulationssignale konzipiert und implementiert werden. 
Ein besonderer Schwerpunkt liegt dabei auf der Untersuchung verschiedener Optimierungsformulierungen, um eine möglichst sparsame Stimulation bei gleichzeitiger effektiver Anfallsunterdrückung zu erreichen. 
Die entwickelten Ansätze sollen abschließend in einer Simulationsumgebung anhand eines biophysikalischen Ganzhirnmodells evaluiert werden.

Die Arbeit umfasst die folgenden Arbeitspakete:
\begin{enumerate}
    \item Einarbeitung in die Themengebiete fokale Epilepsie, elektrische Hirnstimulation und modellprädiktive Regelung sowie Literaturrecherche zu datengetriebenen Vorhersagemodellen neuronaler Aktivität und geeigneten Regelungsansätzen.
    \item Aufsetzen und Konfiguration einer Simulationsumgebung auf Basis eines biophysikalisch motivierten Ganzhirnmodells zur Simulation fokaler epileptischer Anfälle sowie der EEG-Messung und elektrischen Stimulation. 
    \item Konzeption und Implementierung von Vorhersagemodellen für EEG-Signalverläufe sowie Entwicklung modellprädiktiver Regler unter Berücksichtigung sparsamer Stimulationsstrategien.
     \item Evaluation und Vergleich der entwickelten Regelungsansätze anhand von Simulationen hinsichtlich Anfallsunterdrückung, Stimulationsaufwand und Rechenaufwand sowie Erarbeitung von Handlungsempfehlungen für weiterführende Arbeiten. 
\end{enumerate}

\begin{table}[h!]
    \begin{tabular}{l l}
    Tag der Ausgabe:    &   26.05.2026 \\
    Tag der Abgabe:     &   23.11.2026
    \end{tabular}
\end{table}